\subsection{Overview of the finite-particle comparison}
\label{subsec:general-argument-overview}

This appendix expands the proof architecture summarized in
\cref{sec:proof-overview}, especially the
four-step sketch in \cref{subsec:proof-sketch-after-assumptions}.  The population PWGF mechanism itself is treated
as the reference dynamics, in the sense of the Hessian-guided construction of
\citet{yamamoto2025hessianguided}.  The purpose of the present appendix is the
finite-particle empirical error analysis: we quantify how the implementable
particle system tracks the prescribed mean-field PWGF path.

\begin{table}[t]
\centering
\footnotesize
\setlength{\tabcolsep}{3pt}
\renewcommand{\arraystretch}{1.12}
\begin{tabular}{>{\raggedright\arraybackslash}p{0.19\linewidth}
                >{\raggedright\arraybackslash}p{0.22\linewidth}
                >{\raggedright\arraybackslash}p{0.23\linewidth}
                >{\raggedright\arraybackslash}p{0.25\linewidth}}
\toprule
Condition or conclusion & General comparison & Variance--cosine example & Matrix-factorization case study \\
\midrule
Moment structure and weighted smoothness
& Assumed in \cref{assumption-structured-moment-smoothness}
& Verified in \cref{lem:canonical-comparison-assumptions}(i)
& Separate product-model conditions in \cref{assumption:mf-model-restrictions} \\
Gaussian jumps and schedule
& Prescribed by \cref{assumption-common-seed}
& One prescribed common jump, verified in \cref{lem:canonical-comparison-assumptions}(ii)
& --- \\
Moment growth and ODE regularity
& Assumed in \cref{assumption-moment-growth,assumption-uniform-concentration-piecewise-lipschitz}
& Verified in \cref{lem:canonical-comparison-assumptions}(iii)--(iv)
& --- \\
Non-spikiness
& Assumed in \cref{assumption-anti-spikiness-event}
& Verified with $A_{\rm sp}=1$ in \cref{lem:canonical-comparison-assumptions}(v)
& --- \\
Curvature-response mechanism
& Nonlinear stability in \cref{thm:fixed-schedule-high-prob-r-v15}; \(N^\upsilon\)-scaled curvature-response bounds in \cref{cor:fixed-schedule-high-prob-r-rate-v15}
& High-probability response bounds in \cref{lem:canonical-response-control}; \cref{cor:fixed-schedule-high-prob-r-rate-v15} applies with \(\upsilon=\zeta=1/16\)
& Population-level local geometry in \cref{prop:mf-local-geometry}: attraction, transverse descent, and positive spherical product pushforward second variation in two regions \\
Finite-particle conclusion
& Conditional tracking and objective comparison
& Growing-horizon conclusion in \cref{prop:canonical-model-assumptions}
& Open \\
\bottomrule
\end{tabular}
\caption{Relation of the model sections to the comparison theory. The
matrix-factorization column records a separate population curvature-response
case study; --- denotes items outside its scope, and ``Open'' denotes the bridge
to finite-particle tracking.}
\label{tab:assumption-evidence-map}
\end{table}

Classical globally Lipschitz propagation-of-chaos arguments close the
reference-particle stability estimate by Gronwall with a worst-case Lipschitz
constant~\citep{sznitman1991}.
Here the same architecture is driven instead by the curvature propagator
associated with the sampled curvature.

The comparison is organized through the auxiliary particle system
$\bs Y_t=(Y_t^1,\dots,Y_t^N)$:
\[
    \rho_{\bs X_t}
    \quad\longleftrightarrow\quad
    \rho_{\bs Y_t}
    \quad\longleftrightarrow\quad
    \mu_t .
\]
Here $\bs X_t=(X_t^1,\dots,X_t^N)$ is the actual empirical system, whose drift
and jump covariance are evaluated at $\rho_{\bs X_t}$; $\mu_t$ is the population
PWGF path; and $\bs Y_t$ starts from the same initial sample and uses the
population drift and the original population Gaussian field along $\mu_t$.
Procrustes/Wasserstein comparisons of covariance factors provide related
mathematical background \citep{masarotto2019procrustes}. The population-first
conditional construction required here is proved in
\cref{lem:population-first-finite-coupling}: it constructs the empirical
Gaussian coordinates jointly with the field
evaluations, preserving both marginal updates. The reference conditional
product law conditions only on the population fields. Thus the comparison
separates two effects: the pathwise stability of $\bs X_t$ relative to
$\bs Y_t$, and the concentration of the conditionally i.i.d. reference cloud
$\bs Y_t$ around $\mu_t$.

\paragraph{How this appendix matches the main proof sketch.}
The detailed proof follows the same four pieces as
\cref{subsec:proof-sketch-after-assumptions}.
First, \cref{subsec:general-argument-bound-settings} defines the common initial coupling, the common-seed Gaussian
jump construction, the moment-Lipschitz smoothness assumptions, and the
localization event.  Second,
\cref{subsec:general-argument-bound-r_t-ode-continuous} proves the continuous-time contraction-versus-residual bound for
\[
    r_t
    \coloneq
    \frac1{\sqrt N}\norm{\bs X_t-\bs Y_t}_2 .
\]
Third, \cref{subsec:general-argument-r_t-jump-errors} controls the perturbation mismatch at prescribed jump times.  Fourth,
\cref{subsec:general-argument-bound-stochastic-gamma-and-tilde-gamma}, together with
\cref{subsubsec:general-argument-bound-gamma} and
\cref{subsubsec:weighted-factorization-tilde-gamma}, supplies the concentration estimates needed to
close the deterministic inequality.

\paragraph{Core theorem: prescribed common jump schedule.}
The core result of this appendix is the fixed-schedule comparison in
\cref{subsec:fixed-schedule-package}.  In this
layer the population and empirical systems are compared along a prescribed
common jump schedule
\(
0=t_0<t_1<\cdots<t_{\NofJ}\le T
\),

with $\NofJ=O(1)$ in
\Cref{cor:fixed-schedule-high-prob-r-rate-v15}. On the population moment
and sample localization events, and under anti-spikiness, the ODE estimate is
\[
\dot r_t\le-\kappa_t^{Y,N}r_t+c_Dr_t^2+\Gamma_t^Nr_t+\Gamma_t^N.
\]
At each prescribed jump, the comparison uses only pre-jump residuals:
\[
r_{t_\ell}\le(1+a_\ell)r_{t_\ell-}
+a_\ell(\sqrt{\Gamma_{t_\ell-}^N}+\sqrt{\widetilde\Gamma_{t_\ell-}^N}),
\qquad a_\ell=c_J(D)q_p\eta_\ell.
\]
The finite Gaussian coupling is constructed after generating the population
field; the reference product law conditions only on the population fields.
The exact result in \Cref{thm:fixed-schedule-high-prob-r-v15} retains the
nonlinear condition $T(\delta+c_DR)\le1/4$. Its failure probability is bounded
within that nonlinear stability event and the regularity event. The sufficient horizon
$T\le N^{1/4-\upsilon-\zeta}$ in \cref{cor:fixed-schedule-high-prob-r-rate-v15}
also specifies the curvature-response bounds and uniform structural constants it uses.

\paragraph{The two stochastic residuals.}
The continuous residual is
\[
    \Gamma_t^N
    \coloneq
    \sum_{m\ge1}\gamma_m\Delta_{m,t}^N,
    \qquad
    \Delta_{m,t}^N
    \coloneq
    \norm{M_m(\rho_{\bs Y_t})-M_m(\mu_t)} .
\]
It controls the population-to-reference empirical error in the drift, Hessian,
and functional terms.  The jump covariance residual is
\[
    \widetilde\Gamma_t^N
    \coloneq
    \frac1N
    \left\|
    \int
    \mathbf U_t(z)\mathbf U_t(z)^\top
    (\mu_t-\rho_{\bs Y_t})(\dd z)
    \right\|_* ,
\]
where
\[
    \mathbf U_t(z)
    \coloneq
    \bigl(
    \HessW F(\mu_t;Y_t^1,z),\ldots,
    \HessW F(\mu_t;Y_t^N,z)
    \bigr)^\top .
\]
The first residual enters additively along ODE intervals, whereas the second
enters through a square-root covariance perturbation at jumps.  Consequently,
even when both residuals concentrate at nearly $N^{-1/2}$ scale, the jump
contribution generally produces an $N^{-1/4}$ bottleneck.
\Cref{subsubsec:general-argument-bound-gamma} and
\cref{subsubsec:weighted-factorization-tilde-gamma} supply the concentration
estimates for $\Gamma_t^N$ and $\widetilde\Gamma_t^N$, respectively.

\paragraph{From particle tracking to functional error.}
After the fixed-schedule $r_t$ estimate is proved,
\cref{subsec:functional-gap} converts the
tracking bound into a bound on
\[
    \bigl|F(\rho_{\bs X_t})-F(\mu_t)\bigr| .
\]
The key point is that no new stochastic object is introduced: the functional gap
is bounded by the same pair $(r_t,\Gamma_t^N)$. The common weights
\[
 \gamma_m(D)=\mathrm L_{K_{\rm pop},D}a_m
\]
from \cref{lemma-localized-moment-smoothness-consequences} control both
the dynamical comparison and the functional comparison.

\paragraph{An example on a growing horizon.}
\Cref{prop:canonical-model-assumptions} treats a variance-plus-cosine model
with bounded mean-zero initialization and one prescribed Gaussian jump.
Its proof computes the sampled curvature, bounds the probability of
escape, and checks finite-$N$ stability. After positive curvature returns,
the jump curvature-response factor decays, yielding tracking and a bound on the objective
gap to the minimum for $C\log N\le t\le N^{1/8}$.

\paragraph{Fixed-time supplements.}
\Cref{lem:fixed-time-trigger-squared-gap} compares
$\mathcal G_t=\int\|\nabla_WF(\mu_t,x)\|^2\,\mu_t(dx)$ and
$\mathcal G_t^N=N^{-1}\sum_i\|\nabla_WF(\rho_{\mathbf X_t},X_t^i)\|^2$
at one deterministic non-jump time. The margin condition
\eqref{eq:trigger-margin-agreement} gives agreement of threshold decisions
at that time. Separately,
\cref{prop:finite-feature-negative-curvature} transfers negative
population curvature to the reference sample for finite-feature
objectives under a spatial-Hessian condition.
These results are not inputs to the tracking proof and do not control
whole-path curvature-response factors or synchronize adaptive jump times.
